\documentclass[11pt]{article}

\usepackage[margin=1in]{geometry}
\usepackage{amsmath,amssymb,amsthm}
\usepackage{enumitem}
\usepackage{xcolor}
\usepackage{hyperref}

\hypersetup{
  colorlinks=true,
  linkcolor=blue!60!black,
  urlcolor=blue!60!black
}

\newcommand{\R}{\mathbb{R}}
\newcommand{\F}{\mathcal{F}}
\newcommand{\Prob}{\mathbb{P}}
\newcommand{\E}{\mathbb{E}}
\newcommand{\dd}{\,d}
\newcommand{\Del}{\Delta}
\newcommand{\Knowledge}[1]{%
  \par\smallskip\noindent
  \textbf{Knowledge used.}
  {\small\color{blue!60!black}\raggedright #1\par}
  \par\smallskip
}

\title{Chapter 9 Exercises}
\author{\textsc{Solutions to Rick Durrett's}\\
\textit{Probability: Theory and Examples}}
\date{\today}

\begin{document}
\maketitle

\noindent
These solutions use the local Chapter 9 LLM/Viki knowledge set in
\path{Probability/LLM_Viki_Dataset/Chapter9_Knowledge}.  The cited
knowledge IDs refer to entries in
\path{chapter9_knowledge_pieces.jsonl}.

\section*{Section 9.2 Exercises}

\noindent\textbf{Question.}
The local Chapter 9 textbook PDF has no separately numbered exercises
inside Section 9.2, ``Heat Equation.''  The items below solve the
natural verification problems in that section: the martingale
calculation, uniqueness, recovery of the initial condition, and the
kernel-smoothing detail left to the reader in Theorem 9.2.5.

\Knowledge{
\path{durrett_ch9_heat_backward_martingale};
\path{durrett_ch9_heat_equation_solution};
\path{durrett_ch9_heat_kernel_smoothing}.
}

\subsection*{Exercise 9.2.A: the backward heat martingale}

\noindent\textbf{Problem.}
Suppose $u\in C^{1,2}((0,\infty)\times\R^d)$ and
\[
  u_t=\frac12\Del u.
\]
Show that, for fixed $t>0$,
\[
  M_s=u(t-s,B_s),\qquad 0\le s<t,
\]
is a local martingale.

\noindent\textbf{Solution.}
Apply It\^o's formula to the space-time function
$F(s,x)=u(t-s,x)$.  Its derivatives are
\[
  F_s(s,x)=-u_t(t-s,x),\qquad
  \nabla_xF(s,x)=\nabla u(t-s,x),\qquad
  \Del_xF(s,x)=\Del u(t-s,x).
\]
Therefore
\[
  \dd u(t-s,B_s)
  =
  \left[-u_t(t-s,B_s)+\frac12\Del u(t-s,B_s)\right]\dd s
  +\nabla u(t-s,B_s)\cdot\dd B_s.
\]
The drift vanishes because $u_t=(1/2)\Del u$.  Thus
\[
  M_s-M_0
  =
  \int_0^s \nabla u(t-r,B_r)\cdot\dd B_r.
\]
The stochastic integral is a local martingale after the usual
localization, for instance by stopping when
\[
  \int_0^s |\nabla u(t-r,B_r)|^2\dd r
\]
exceeds $n$.  Hence $M_s$ is a local martingale on $[0,t)$.

\subsection*{Exercise 9.2.B: uniqueness of bounded heat solutions}

\noindent\textbf{Problem.}
Let $u$ be a bounded solution of
\[
  u_t=\frac12\Del u,\qquad
  u(0,x)=f(x),
\]
where $f$ is bounded and continuous.  Prove that
\[
  u(t,x)=\E_x f(B_t).
\]

\noindent\textbf{Solution.}
Fix $t>0$ and start Brownian motion at $B_0=x$.  By Exercise 9.2.A,
$M_s=u(t-s,B_s)$ is a local martingale for $0\le s<t$.  Since $u$ is
bounded, the localized martingales are uniformly bounded; hence
$M_s$ is an honest bounded martingale on each interval $[0,t-\epsilon]$,
and the martingale convergence theorem gives
\[
  M_t:=\lim_{s\uparrow t}M_s
\]
almost surely and in $L^1$.

Because $B_s\to B_t$ almost surely and $u$ is continuous at
$(0,B_t)$ with boundary value $f(B_t)$,
\[
  M_t
  =
  \lim_{s\uparrow t}u(t-s,B_s)
  =
  f(B_t)
  \qquad \text{a.s.}
\]
Taking expectations in the bounded martingale gives
\[
  u(t,x)=M_0=\E_x M_t=\E_x f(B_t).
\]
Thus a bounded solution is unique and must be the Brownian heat
semigroup applied to $f$.

\subsection*{Exercise 9.2.C: the initial condition for the Brownian formula}

\noindent\textbf{Problem.}
Let
\[
  v(t,x)=\E_x f(B_t)=\E f(x+B_t),
\]
where $f$ is bounded and continuous.  Show that
\[
  v(t_n,x_n)\to f(x)
\]
whenever $t_n\downarrow0$ and $x_n\to x$.

\noindent\textbf{Solution.}
Write $B_{t_n}\stackrel{d}{=}\sqrt{t_n}Z$, where
$Z$ is a standard $d$-dimensional normal vector.  Then
\[
  v(t_n,x_n)=\E f(x_n+\sqrt{t_n}Z).
\]
For each fixed $Z$, $x_n+\sqrt{t_n}Z\to x$.  Since $f$ is continuous,
\[
  f(x_n+\sqrt{t_n}Z)\to f(x)
  \qquad \text{a.s.}
\]
Also $|f(x_n+\sqrt{t_n}Z)|\le \|f\|_\infty$.  The bounded convergence
theorem gives
\[
  v(t_n,x_n)\to \E f(x)=f(x).
\]
So the Brownian formula has the required initial trace.

\subsection*{Exercise 9.2.D: heat-kernel smoothing and the PDE}

\noindent\textbf{Problem.}
Fill in the details in Theorem 9.2.5.  For bounded continuous $f$, prove
that
\[
  v(t,x)=\E_x f(B_t)
  =
  \int_{\R^d}p_t(x-y)f(y)\dd y,
  \qquad
  p_t(z)=(2\pi t)^{-d/2}e^{-|z|^2/(2t)},
\]
belongs to $C^{1,2}((0,\infty)\times\R^d)$ and satisfies
\[
  v_t=\frac12\Del v.
\]

\noindent\textbf{Solution.}
Fix $0<a<b<\infty$ and a compact set $K\subset\R^d$.  We first justify
differentiating under the integral for $(t,x)\in[a,b]\times K$.
For a multi-index $\alpha$ and an integer $m\ge0$, the derivatives
\[
  \partial_t^m\partial_x^\alpha p_t(x-y)
\]
are finite sums of terms of the form
\[
  C\,t^{-q}\,(x-y)^\beta p_t(x-y),
\]
where $q$ and $\beta$ depend on $m,\alpha,d$.  On $t\in[a,b]$ these are
bounded by an integrable Gaussian-tail majorant
\[
  C_{a,b,K}\,(1+|y|^N)e^{-c|y|^2}
\]
for suitable constants $C_{a,b,K},N,c>0$.  Since $f$ is bounded, the
dominated convergence theorem allows us to differentiate under the
integral for one time derivative and for two spatial derivatives.
Thus $v\in C^{1,2}$ on $(0,\infty)\times\R^d$.

It remains to identify the PDE.  The Gaussian kernel satisfies
\[
  \partial_t p_t(z)
  =
  \left(-\frac d{2t}+\frac{|z|^2}{2t^2}\right)p_t(z),
\]
while
\[
  \partial_{z_i}p_t(z)=-\frac{z_i}{t}p_t(z),
  \qquad
  \partial_{z_i z_i}p_t(z)
  =
  \left(\frac{z_i^2}{t^2}-\frac1t\right)p_t(z).
\]
Summing over $i$ gives
\[
  \Del p_t(z)
  =
  \left(\frac{|z|^2}{t^2}-\frac d t\right)p_t(z).
\]
Hence
\[
  \partial_t p_t(z)=\frac12\Del p_t(z).
\]
Differentiating the convolution under the integral,
\[
  v_t(t,x)
  =
  \int_{\R^d}\partial_t p_t(x-y)f(y)\dd y
  =
  \frac12\int_{\R^d}\Del_x p_t(x-y)f(y)\dd y
  =
  \frac12\Del v(t,x).
\]
Therefore $v$ is a classical solution of the heat equation for $t>0$.
Together with Exercise 9.2.C, this proves existence of a bounded
solution with initial value $f$; Exercise 9.2.B proves uniqueness.

\section*{Section 9.3 Exercises}

\noindent\textbf{Question.}
The local Chapter 9 textbook PDF also has no separately numbered
exercise block inside Section 9.3, ``Inhomogeneous Heat Equation.''
The solutions below spell out the section's verification problems for
the equation
\[
  u_t=\frac12\Del u+g,
  \qquad
  u(0,x)=0,
\]
with bounded continuous $g$.  For nonzero initial condition $f$, add
the homogeneous heat solution $\E_x f(B_t)$ from Section 9.2.

\Knowledge{
\path{durrett_ch9_duhamel_inhomogeneous_heat};
\path{durrett_ch9_inhomogeneous_heat_regularity};
\path{durrett_ch9_heat_kernel_smoothing}.
}

\subsection*{Exercise 9.3.A: the source-term martingale}

\noindent\textbf{Problem.}
Suppose $u\in C^{1,2}$ and
\[
  u_t=\frac12\Del u+g.
\]
Show that, for fixed $t>0$,
\[
  M_s=u(t-s,B_s)+\int_0^s g(B_r)\dd r,\qquad 0\le s<t,
\]
is a local martingale.

\noindent\textbf{Solution.}
As in Section 9.2, apply It\^o's formula to $u(t-s,B_s)$:
\[
  \dd u(t-s,B_s)
  =
  \left[-u_t(t-s,B_s)+\frac12\Del u(t-s,B_s)\right]\dd s
  +\nabla u(t-s,B_s)\cdot\dd B_s.
\]
Since $u_t=(1/2)\Del u+g$,
\[
  -u_t+\frac12\Del u=-g.
\]
Hence
\[
  \dd\left(u(t-s,B_s)+\int_0^s g(B_r)\dd r\right)
  =
  \nabla u(t-s,B_s)\cdot\dd B_s.
\]
After localizing the stochastic integral by its quadratic variation,
the right side is a martingale.  Therefore $M_s$ is a local martingale.

\subsection*{Exercise 9.3.B: uniqueness and Duhamel's formula}

\noindent\textbf{Problem.}
Assume $u$ is bounded on $[0,T]\times\R^d$ for each $T<\infty$,
satisfies
\[
  u_t=\frac12\Del u+g,\qquad u(0,x)=0,
\]
and $g$ is bounded.  Prove that
\[
  u(t,x)=\E_x\int_0^t g(B_s)\dd s.
\]

\noindent\textbf{Solution.}
Fix $t>0$ and $B_0=x$.  By Exercise 9.3.A,
\[
  M_s=u(t-s,B_s)+\int_0^s g(B_r)\dd r,\qquad 0\le s<t,
\]
is a local martingale.  Since $u$ and $g$ are bounded on the relevant
finite time interval, $M_s$ is bounded by a deterministic constant, so
the localized martingale is uniformly integrable and $M_s$ is an honest
martingale.

Letting $s\uparrow t$, continuity at time $0$ gives
\[
  \lim_{s\uparrow t}u(t-s,B_s)=u(0,B_t)=0,
\]
and bounded convergence gives
\[
  M_t=\int_0^t g(B_r)\dd r.
\]
Taking expectations,
\[
  u(t,x)=M_0=\E_xM_t
  =
  \E_x\int_0^t g(B_s)\dd s.
\]
This proves uniqueness among bounded solutions and identifies the
solution as Duhamel's formula.

\subsection*{Exercise 9.3.C: smooth formula implies the PDE}

\noindent\textbf{Problem.}
Let
\[
  v(t,x)=\E_x\int_0^t g(B_s)\dd s.
\]
Assume $v\in C^{1,2}((0,\infty)\times\R^d)$.  Show that
\[
  v_t=\frac12\Del v+g.
\]

\noindent\textbf{Solution.}
Use the Markov property to write, for $h>0$,
\[
  v(t+h,x)
  =
  \E_x\int_0^{t+h}g(B_s)\dd s
  =
  \E_x\int_0^h g(B_s)\dd s+\E_x v(t,B_h).
\]
Subtract $v(t,x)$:
\[
  v(t+h,x)-v(t,x)
  =
  \E_x\{v(t,B_h)-v(t,x)\}
  +\E_x\int_0^h g(B_s)\dd s.
\]
Because $v$ is $C^{1,2}$ in space, Taylor's formula at $x$ gives
\[
  \E_x\{v(t,B_h)-v(t,x)\}
  =
  \frac12\Del v(t,x)h+o(h),
\]
using $\E_x(B_h^i-x_i)=0$,
$\E_x[(B_h^i-x_i)(B_h^j-x_j)]=h\delta_{ij}$.  Since $g$ is continuous
and $B_s\to x$ in probability as $s\downarrow0$,
\[
  \frac1h\E_x\int_0^h g(B_s)\dd s\to g(x).
\]
Divide by $h$ and let $h\downarrow0$:
\[
  v_t(t,x)=\frac12\Del v(t,x)+g(x).
\]

\subsection*{Exercise 9.3.D: initial condition}

\noindent\textbf{Problem.}
Show that
\[
  v(t,x)=\E_x\int_0^t g(B_s)\dd s
\]
satisfies $v(t,x)\to0$ as $t\downarrow0$, uniformly in $x$, when
$g$ is bounded.

\noindent\textbf{Solution.}
If $\|g\|_\infty\le M$, then
\[
  |v(t,x)|
  \le
  \E_x\int_0^t |g(B_s)|\dd s
  \le Mt.
\]
Therefore $\sup_x |v(t,x)|\le Mt\to0$.  This is exactly the initial
condition $v(0,x)=0$.

\subsection*{Exercise 9.3.E: kernel form and regularity}

\noindent\textbf{Problem.}
Write the Duhamel solution in heat-kernel form and explain why
Holder continuity of $g$ is a natural sufficient condition for
$v\in C^{1,2}$.

\noindent\textbf{Solution.}
By Fubini's theorem and the Brownian transition density,
\[
  v(t,x)
  =
  \int_0^t \E_x g(B_s)\dd s
  =
  \int_0^t\int_{\R^d}p_s(x-y)g(y)\dd y\dd s.
\]
Equivalently,
\[
  v(t,x)=\int_{\R^d}G_t(x,y)g(y)\dd y,
  \qquad
  G_t(x,y)=\int_0^t p_s(x-y)\dd s.
\]
Formally differentiating under the integral is delicate near $s=0$,
because derivatives of $p_s$ become singular as $s\downarrow0$ and
$y\to x$.  The cancellation supplied by
\[
  g(y)=g(x)+O(|y-x|^\alpha)
\]
when $g$ is Holder continuous offsets this singularity.  More
concretely, split
\[
  P_sg(x)-g(x)
  =
  \int_{\R^d}p_s(z)\{g(x+z)-g(x)\}\dd z.
\]
If $g$ is Holder with exponent $\alpha>0$, then
\[
  |P_sg(x)-g(x)|
  \le
  C\E|B_s|^\alpha
  =
  C' s^{\alpha/2}.
\]
This is the small-time gain that makes the singular kernel estimates
integrable in the parabolic regularity argument.  The standard
Schauder estimate for the heat equation then gives
\[
  v\in C^{1,2}((0,\infty)\times\R^d),
\]
and Exercise 9.3.C proves that this classical function solves
\[
  v_t=\frac12\Del v+g,\qquad v(0,x)=0.
\]
Thus Holder continuity is not just a cosmetic assumption; it is the
condition that lets the source term smooth into a classical solution.

\section*{Section 9.4 Exercises}

\noindent\textbf{Question.}
The local Chapter 9 textbook PDF has no separately numbered exercise
block inside Section 9.4, ``Feynman-Kac Formula.''  The solutions below
work through the section's verification problems for
\[
  u_t=\frac12\Del u+c(x)u,
  \qquad
  u(0,x)=f(x),
\]
where $c$ and $f$ are bounded and continuous.

\Knowledge{
\path{durrett_ch9_feynman_kac_forward};
\path{durrett_ch9_feynman_kac_smoothness};
\path{durrett_ch9_heat_equation_solution};
\path{durrett_ch9_duhamel_inhomogeneous_heat}.
}

\subsection*{Exercise 9.4.A: the weighted martingale}

\noindent\textbf{Problem.}
Let
\[
  A_s=\int_0^s c(B_r)\dd r.
\]
Suppose $u\in C^{1,2}$ and
\[
  u_t=\frac12\Del u+c(x)u.
\]
Show that, for fixed $t>0$,
\[
  M_s=u(t-s,B_s)e^{A_s},\qquad 0\le s<t,
\]
is a local martingale.

\noindent\textbf{Solution.}
The process $A_s$ has finite variation, so the ordinary product rule
and It\^o's formula give
\[
  \dd\{u(t-s,B_s)e^{A_s}\}
  =
  e^{A_s}\dd u(t-s,B_s)
  +
  u(t-s,B_s)e^{A_s}\dd A_s.
\]
There is no quadratic covariation term because $e^{A_s}$ has finite
variation.  Also
\[
  \dd A_s=c(B_s)\dd s
\]
and
\[
  \dd u(t-s,B_s)
  =
  \left[-u_t(t-s,B_s)+\frac12\Del u(t-s,B_s)\right]\dd s
  +
  \nabla u(t-s,B_s)\cdot\dd B_s.
\]
Therefore
\[
  \dd M_s
  =
  e^{A_s}
  \left[-u_t+\frac12\Del u+cu\right](t-s,B_s)\dd s
  +
  e^{A_s}\nabla u(t-s,B_s)\cdot\dd B_s.
\]
Since $u_t=(1/2)\Del u+cu$, the drift vanishes.  After localizing the
stochastic integral by its quadratic variation,
\[
  M_s-M_0
  =
  \int_0^s e^{A_r}\nabla u(t-r,B_r)\cdot\dd B_r
\]
is a local martingale.  Hence $M_s$ is a local martingale.

\subsection*{Exercise 9.4.B: uniqueness and the Feynman-Kac formula}

\noindent\textbf{Problem.}
Assume $u$ is bounded on $[0,T]\times\R^d$ for each $T<\infty$ and
satisfies
\[
  u_t=\frac12\Del u+c(x)u,\qquad u(0,x)=f(x),
\]
with $c$ bounded.  Prove that
\[
  u(t,x)
  =
  \E_x\left[
    f(B_t)\exp\left(\int_0^t c(B_s)\dd s\right)
  \right].
\]

\noindent\textbf{Solution.}
Fix $t>0$ and start Brownian motion at $B_0=x$.  By Exercise 9.4.A,
\[
  M_s=u(t-s,B_s)\exp\left(\int_0^s c(B_r)\dd r\right)
\]
is a local martingale on $[0,t)$.  If $|c|\le C$ and $u$ is bounded on
$[0,t]\times\R^d$, then
\[
  |M_s|\le e^{Ct}\sup_{0\le r\le t,y\in\R^d}|u(r,y)|.
\]
Thus $M_s$ is bounded and therefore an honest uniformly integrable
martingale.

Let $s\uparrow t$.  Since $B_s\to B_t$ almost surely and $u$ is
continuous at the initial boundary,
\[
  \lim_{s\uparrow t}u(t-s,B_s)=u(0,B_t)=f(B_t).
\]
Also $\int_0^s c(B_r)\dd r\to\int_0^t c(B_r)\dd r$.  Hence
\[
  M_t
  =
  f(B_t)\exp\left(\int_0^t c(B_r)\dd r\right).
\]
Taking expectations in the martingale identity $M_0=\E_xM_t$ gives the
claimed formula.  In particular, this proves uniqueness among bounded
solutions.

\subsection*{Exercise 9.4.C: a smooth Feynman-Kac formula solves the PDE}

\noindent\textbf{Problem.}
Let
\[
  v(t,x)=
  \E_x\left[
    f(B_t)\exp\left(\int_0^t c(B_s)\dd s\right)
  \right].
\]
Assume $v\in C^{1,2}((0,\infty)\times\R^d)$.  Show that
\[
  v_t=\frac12\Del v+c(x)v.
\]

\noindent\textbf{Solution.}
The Markov property gives, for $h>0$,
\[
  v(t+h,x)
  =
  \E_x\left[
    e^{\int_0^h c(B_s)\dd s}v(t,B_h)
  \right].
\]
Subtract $v(t,x)$ and add-subtract $v(t,B_h)$:
\[
  v(t+h,x)-v(t,x)
  =
  \E_x\{v(t,B_h)-v(t,x)\}
  +
  \E_x\left[
    v(t,B_h)\left(e^{\int_0^h c(B_s)\dd s}-1\right)
  \right].
\]
Since $v$ is $C^{1,2}$ in space,
\[
  \E_x\{v(t,B_h)-v(t,x)\}
  =
  \frac12\Del v(t,x)h+o(h).
\]
For the second term, boundedness of $c$ gives
\[
  e^{\int_0^h c(B_s)\dd s}-1
  =
  \int_0^h c(B_r)\exp\left(\int_r^h c(B_s)\dd s\right)\dd r.
\]
As $h\downarrow0$, $B_r\to x$ in probability uniformly on small time
intervals and $v(t,B_h)\to v(t,x)$ in probability.  The integrand is
bounded on compact time intervals, so dominated convergence yields
\[
  \frac1h
  \E_x\left[
    v(t,B_h)\left(e^{\int_0^h c(B_s)\dd s}-1\right)
  \right]
  \to c(x)v(t,x).
\]
Divide by $h$ and let $h\downarrow0$ to obtain
\[
  v_t(t,x)=\frac12\Del v(t,x)+c(x)v(t,x).
\]

\subsection*{Exercise 9.4.D: initial condition}

\noindent\textbf{Problem.}
Show that the Feynman-Kac formula satisfies the initial condition:
\[
  v(t_n,x_n)\to f(x)
\]
whenever $t_n\downarrow0$ and $x_n\to x$.

\noindent\textbf{Solution.}
Write
\[
  A_t=\int_0^t c(B_s)\dd s.
\]
If $|c|\le C$, then $|A_t|\le Ct$, so $e^{A_t}\to1$ uniformly as
$t\downarrow0$.  Therefore
\[
  |v(t,x)-\E_xf(B_t)|
  \le
  \|f\|_\infty \E_x|e^{A_t}-1|
  \le
  \|f\|_\infty(e^{Ct}-1)\to0.
\]
By the Section 9.2 initial-condition result,
\[
  \E_{x_n}f(B_{t_n})\to f(x).
\]
Combining the two estimates gives
\[
  v(t_n,x_n)\to f(x).
\]

\subsection*{Exercise 9.4.E: reducing smoothness to Sections 9.2 and 9.3}

\noindent\textbf{Problem.}
Assume $f$ and $c$ are Holder continuous and bounded.  Explain the
identity used in Durrett's proof of Theorem 9.4.5 and why it reduces
smoothness of the Feynman-Kac formula to the heat and inhomogeneous
heat results.

\noindent\textbf{Solution.}
Let
\[
  A_t=\int_0^t c(B_r)\dd r.
\]
Since $A_t$ is continuous and of finite variation,
\[
  e^{A_t}-1
  =
  \int_0^t c(B_s)e^{A_t-A_s}\dd s
  =
  \int_0^t c(B_s)
      \exp\left(\int_s^t c(B_r)\dd r\right)\dd s.
\]
Multiplying by $f(B_t)$ and taking expectations,
\[
  v(t,x)
  =
  \E_xf(B_t)
  +
  \int_0^t
    \E_x\left[
      c(B_s)
      \exp\left(\int_s^t c(B_r)\dd r\right)f(B_t)
    \right]\dd s.
\]
Conditioning on $\F_s$ and using the Markov property turns the second
integrand into
\[
  \E_x\left[c(B_s)v(t-s,B_s)\right].
\]
Thus
\[
  v(t,x)=v_1(t,x)+v_2(t,x),
\]
where
\[
  v_1(t,x)=\E_xf(B_t),
  \qquad
  v_2(t,x)=\int_0^t\E_x\{c(B_s)v(t-s,B_s)\}\dd s.
\]
The first term is the homogeneous heat solution from Section 9.2, so
$v_1\in C^{1,2}$ for $t>0$.  The second term has the form of the
inhomogeneous heat solution from Section 9.3 with source
\[
  g(r,x)=c(x)v(r,x).
\]
Boundedness of $c$ and $f$ makes $v$ locally bounded in time.  Holder
continuity of $f$ gives local Holder control of $v_1$ in $x$, while the
inhomogeneous heat estimates give the same kind of local Holder control
for $v_2$.  Since $c$ is Holder, the product $g(r,x)=c(x)v(r,x)$ is
locally Holder in the spatial variable.  Section 9.3's regularity
result then gives $v_2\in C^{1,2}$, and hence
\[
  v=v_1+v_2\in C^{1,2}.
\]
Exercise 9.4.C then shows that this smooth Feynman-Kac formula solves
the PDE.

\section*{Section 9.5 Exercises}

\Knowledge{
\path{durrett_ch9_dirichlet_problem_representation};
\path{durrett_ch9_mean_value_harmonicity};
\path{durrett_ch9_interior_smoothness_harmonic};
\path{durrett_ch9_regular_boundary_point};
\path{durrett_ch9_boundary_continuity_regular_points};
\path{durrett_ch9_cone_condition}.
}

\subsection*{Exercise 9.5.1}

\noindent\textbf{Question.}
Suppose $h$ is bounded and harmonic on $\R^d$.  Prove that $h$ is
constant.

\noindent\textbf{Solution.}
Let $B_t$ be Brownian motion started at $x$.  Since $h$ is harmonic,
$h(B_{t\wedge S_R})$ is a bounded martingale, where
$S_R=\inf\{t:|B_t-x|=R\}$.  Letting $R\to\infty$ and using bounded
convergence gives
\[
  h(x)=\E h(x+B_t),\qquad t>0.
\]
Thus, for any $x,y\in\R^d$,
\[
  h(x)-h(y)
  =
  \int_{\R^d} h(z)\{p_t(z-x)-p_t(z-y)\}\dd z.
\]
If $\|h\|_\infty\le M$, then
\[
  |h(x)-h(y)|
  \le
  M\int_{\R^d}|p_t(z-x)-p_t(z-y)|\dd z.
\]
The two Gaussian densities have the same covariance matrix $tI$ and
means separated by the fixed vector $x-y$.  As $t\to\infty$, their
total variation distance tends to $0$.  Hence $|h(x)-h(y)|=0$.
Since $x$ and $y$ were arbitrary, $h$ is constant.

\subsection*{Exercise 9.5.2}

\noindent\textbf{Question.}
Suppose $f\in C^2$ is superharmonic, takes values in $[0,\infty]$, and
is continuous.  Show that in dimensions $1$ and $2$, $f$ must be
constant.  Give an example of the form $f(x)=g(|x|)$ in $d\ge3$.

\noindent\textbf{Solution.}
Use the convention that superharmonic means $\Del f\le0$.  It\^o's
formula gives
\[
  f(B_t)-f(B_0)
  =
  \int_0^t \nabla f(B_s)\cdot\dd B_s
  +
  \frac12\int_0^t \Del f(B_s)\dd s.
\]
After localization, $f(B_t)$ is therefore a nonnegative local
supermartingale, hence a supermartingale.

Let $x,y\in\R^d$ and let
\[
  T_\varepsilon=\inf\{t:|B_t-y|\le\varepsilon\}
\]
for Brownian motion started at $x$.  In dimensions $1$ and $2$,
Brownian motion is recurrent, so $T_\varepsilon<\infty$ almost surely.
Optional stopping at $T_\varepsilon\wedge n$ and nonnegativity give
\[
  f(x)\ge \E_x f(B_{T_\varepsilon}); 
\]
then let $n\to\infty$.  Since $B_{T_\varepsilon}$ lies on the sphere
of radius $\varepsilon$ about $y$, continuity implies
\[
  \lim_{\varepsilon\downarrow0}\E_x f(B_{T_\varepsilon})=f(y).
\]
Thus $f(x)\ge f(y)$.  Reversing $x$ and $y$ gives $f(y)\ge f(x)$, so
$f$ is constant.

For $d\ge3$, a smooth nonconstant radial example is
\[
  f(x)=g(|x|)
  =
  (1+|x|^2)^{-(d-2)/2}.
\]
Let $\alpha=(d-2)/2$ and $r=|x|$.  For $g(r)=(1+r^2)^{-\alpha}$,
\[
  g''(r)+\frac{d-1}{r}g'(r)
  =
  2\alpha(1+r^2)^{-\alpha-2}
  \{-d+(2\alpha+2-d)r^2\}.
\]
Since $2\alpha+2-d=0$, this becomes
\[
  \Del f(x)
  =
  -d(d-2)(1+|x|^2)^{-(d+2)/2}\le0.
\]
So $f$ is nonnegative, smooth, radial, superharmonic, and nonconstant.

\subsection*{Exercise 9.5.3}

\noindent\textbf{Question.}
Let $\tau=\inf\{t:B_t\notin G\}$.  Suppose $G$ is connected and open,
all boundary points are regular, and $\Prob_x(\tau<\infty)<1$ in $G$.
Show that all bounded solutions of the Dirichlet problem with boundary
value $0$ on $\partial G$ have the form
\[
  u(x)=c\Prob_x(\tau=\infty).
\]

\noindent\textbf{Solution.}
Put
\[
  q(x)=\Prob_x(\tau=\infty).
\]
By the strong Markov property, $q$ is harmonic in $G$.  Since every
boundary point is regular, $q(x)\to0$ as $x$ approaches $\partial G$.
Thus $q$ itself is a bounded solution of the zero-boundary Dirichlet
problem.  The assumption $\Prob_x(\tau<\infty)<1$ says $q(x)>0$ in $G$.

Let $u$ be another bounded solution with zero boundary data and let
$M=\|u\|_\infty$.  Stopping $u(B_t)$ at $\tau$ and using the regular
zero boundary condition gives
\[
  u(x)=\E_x\{u(B_t);\tau>t\}.
\]
Therefore
\[
  |u(x)|\le M\Prob_x(\tau>t).
\]
Letting $t\to\infty$ gives
\[
  |u(x)|\le Mq(x).
\]
Hence $h(x)=u(x)/q(x)$ is bounded in $G$.

Now use the Doob $q$-transform, i.e. Brownian motion conditioned never
to leave $G$:
\[
  \Prob_x^q(A)
  =
  \frac{1}{q(x)}
  \E_x\{\mathbf{1}_A q(B_t);t<\tau\},
  \qquad A\in\mathcal{F}_t.
\]
The ratio $h=u/q$ is harmonic for this conditioned process.  Since $G$
is connected and all finite boundary points are regular, the only
remaining boundary point for the conditioned process is the ideal point
``at infinity.''  The bounded martingale $h(B_t)$ therefore has a
constant terminal value under the conditioned law.  Calling this
constant $c$, we get
\[
  h(x)=\E_x^q h(B_t)=c,
\]
and hence
\[
  u(x)=cq(x)=c\Prob_x(\tau=\infty).
\]

Equivalently, this is the maximum-principle content of the hint:
adding a large multiple of $q$ to $\pm u$ produces nonnegative
zero-boundary harmonic functions, and the only possible freedom is the
amount of mass assigned to escape through infinity.

\subsection*{Exercise 9.5.4}

\noindent\textbf{Question.}
Let $\tau=\inf\{t:B_t\notin G\}$ and consider the Dirichlet problem in
a connected open set $G$ in $d\ge3$, where all boundary points are
regular and $\Prob_x(\tau<\infty)<1$ in $G$.  The extracted statement
again has boundary condition $f\equiv0$.  Conclude that all bounded
solutions have the form
\[
  u(x)=c\Prob_x(\tau=\infty).
\]

\noindent\textbf{Solution.}
For the zero-boundary statement, this is exactly Exercise 9.5.3.  The
dimension assumption $d\ge3$ is the natural transient setting in which
escape to infinity may have positive probability.

The usual general form is worth recording.  If the boundary value is a
bounded continuous function $f$, one particular solution is
\[
  v(x)=\E_x\{f(B_\tau);\tau<\infty\}.
\]
If $u$ is any other bounded solution with the same boundary data, then
$w=u-v$ is a bounded solution with zero boundary data.  By Exercise
9.5.3,
\[
  w(x)=c\Prob_x(\tau=\infty).
\]
Hence all bounded solutions with boundary value $f$ are
\[
  u(x)=
  \E_x\{f(B_\tau);\tau<\infty\}
  +
  c\Prob_x(\tau=\infty).
\]
When $f\equiv0$, this reduces to the displayed formula in the question.

\subsection*{Exercise 9.5.5}

\noindent\textbf{Question.}
Define a flat cone $\overline V(y,v,a)$ to be the intersection of
$V(y,v,a)$ with a $(d-1)$-dimensional hyperplane that contains the line
$\{y+\theta v:\theta\in\R\}$.  Show that Theorem 9.5.10 remains true if
``cone'' is replaced by ``flat cone.''

\noindent\textbf{Solution.}
By translating and rotating, take $y=0$ and take the hyperplane to be
\[
  H=\{x\in\R^d:x_d=0\}.
\]
Write Brownian motion as $B_t=(X_t,Y_t)$, where $X_t$ is
$(d-1)$-dimensional Brownian motion in $H$ and $Y_t$ is an independent
one-dimensional Brownian motion perpendicular to $H$.

The zero set of $Y_t$ accumulates at $0$ almost surely.  At those times
Brownian motion lies in the hyperplane $H$.  The trace of $B_t$ on
$H$, parametrized by the local time of $Y_t$ at $0$, is a rotationally
symmetric Cauchy process in $H$.  Starting from $0$, this trace process
enters every relatively open cone in $H$ immediately with probability
one.  Therefore, if $G^c$ contains a flat cone with vertex $y$, then
Brownian motion started at $y$ hits $G^c$ at arbitrarily small positive
times almost surely.

Thus
\[
  \Prob_y(\tau_G=0)=1,
\]
so $y$ is a regular boundary point.  This is exactly the conclusion of
Theorem 9.5.10 with ``cone'' replaced by ``flat cone.''

\section*{Section 9.6 Exercises}

\noindent\textbf{Question.}
The local Chapter 9 PDF has no separately numbered exercise block in
Section 9.6, ``Green's Functions and Potential Kernels.''  The items
below solve the section's natural verification problems: occupation
times of balls, the potential-kernel computations in dimensions
$d\ge3$, $d=2$, and $d=1$, and the normalization
$(1/2)\Del G=-\delta_0$.

\Knowledge{
\path{durrett_ch9_newtonian_potential_kernel};
\path{durrett_ch9_green_function_killed_domain};
\path{durrett_ch9_radial_harmonic_functions}.
}

\subsection*{Exercise 9.6.A: occupation time of a ball}

\noindent\textbf{Problem.}
Let $D=B(0,r)$.  Show that
\[
  \Prob_x\left(\int_0^\infty \mathbf{1}_D(B_t)\dd t=\infty\right)=1
  \quad\text{in } d\le2,
\]
while
\[
  \E_x\int_0^\infty \mathbf{1}_D(B_t)\dd t<\infty
  \quad\text{in } d\ge3.
\]

\noindent\textbf{Solution.}
For $d\le2$, Brownian motion is recurrent.  Thus it enters $D$
infinitely often almost surely.  Let $G=B(0,2r)$ and define successive
entrance and exit times
\[
  S_k=\inf\{t>T_{k-1}:B_t\in D\},
  \qquad
  T_k=\inf\{t>S_k:B_t\notin G\},
\]
with $T_0=0$.  By the strong Markov property, the occupation increments
\[
  X_k=\int_{S_k}^{T_k}\mathbf{1}_D(B_t)\dd t
\]
are independent after conditioning on the entrance point, and each has
a positive mean bounded away from $0$ because Brownian motion started
in the compact ball $D$ spends positive expected time in $D$ before
leaving $G$.  Hence, by the strong law applied along visits,
\[
  \sum_{k=1}^\infty X_k=\infty
  \qquad\text{a.s.}
\]
Therefore the total occupation time of $D$ is infinite almost surely.

For $d\ge3$, use Fubini and the transition density
\[
  p_t(x,y)=(2\pi t)^{-d/2}e^{-|x-y|^2/(2t)}.
\]
For nonnegative $f$,
\[
  \E_x\int_0^\infty f(B_t)\dd t
  =
  \int_{\R^d}\left(\int_0^\infty p_t(x,y)\dd t\right)f(y)\dd y.
\]
The inner integral is finite for $x\ne y$ when $d\ge3$ and equals a
constant times $|x-y|^{2-d}$, as computed in Exercise 9.6.B.  Taking
$f=\mathbf{1}_D$, the singularity is locally integrable:
\[
  \int_{B(0,r)} |x-y|^{2-d}\dd y<\infty,
\]
because near $y=x$ the polar-coordinate integrand is
$\rho^{2-d}\rho^{d-1}=\rho$.  Thus
\[
  \E_x\int_0^\infty \mathbf{1}_D(B_t)\dd t<\infty.
\]

\subsection*{Exercise 9.6.B: the potential kernel in $d\ge3$}

\noindent\textbf{Problem.}
Compute
\[
  G(x,y)=\int_0^\infty p_t(x,y)\dd t
\]
for $d\ge3$.

\noindent\textbf{Solution.}
Let $a=|x-y|$.  For $a>0$,
\[
  G(x,y)=
  \int_0^\infty (2\pi t)^{-d/2}e^{-a^2/(2t)}\dd t.
\]
Set $s=a^2/(2t)$, so $t=a^2/(2s)$ and
\[
  \dd t=-\frac{a^2}{2s^2}\dd s.
\]
Then
\[
\begin{aligned}
  G(x,y)
  &=
  \int_\infty^0
  \left(2\pi\frac{a^2}{2s}\right)^{-d/2}
  e^{-s}
  \left(-\frac{a^2}{2s^2}\right)\dd s  \\
  &=
  \frac{a^{2-d}}{2\pi^{d/2}}
  \int_0^\infty s^{d/2-2}e^{-s}\dd s  \\
  &=
  \frac{\Gamma(d/2-1)}{2\pi^{d/2}}|x-y|^{2-d}.
\end{aligned}
\]
The gamma integral is finite exactly when $d/2-1>0$, i.e. $d\ge3$.

\subsection*{Exercise 9.6.C: the modified kernel in $d=1$}

\noindent\textbf{Problem.}
In dimension $1$, define
\[
  G(x,y)=\int_0^\infty\{p_t(x,y)-p_t(0,0)\}\dd t.
\]
Show that
\[
  G(x,y)=-|x-y|.
\]

\noindent\textbf{Solution.}
Let $a=|x-y|$.  Then
\[
  G(x,y)
  =
  \frac1{\sqrt{2\pi}}\int_0^\infty
  \left(e^{-a^2/(2t)}-1\right)t^{-1/2}\dd t.
\]
With $u=a^2/(2t)$,
\[
  G(x,y)
  =
  -\frac{a}{2\sqrt{\pi}}\int_0^\infty (1-e^{-u})u^{-3/2}\dd u.
\]
Use
\[
  1-e^{-u}=\int_0^u e^{-s}\dd s
\]
and Fubini:
\[
\begin{aligned}
  \int_0^\infty (1-e^{-u})u^{-3/2}\dd u
  &=
  \int_0^\infty e^{-s}\int_s^\infty u^{-3/2}\dd u\,\dd s  \\
  &=
  2\int_0^\infty e^{-s}s^{-1/2}\dd s
  =
  2\sqrt{\pi}.
\end{aligned}
\]
Therefore
\[
  G(x,y)=-a=-|x-y|.
\]

\subsection*{Exercise 9.6.D: the modified kernel in $d=2$}

\noindent\textbf{Problem.}
In dimension $2$, define
\[
  G(x,y)=\int_0^\infty\{p_t(x,y)-p_t(0,e_1)\}\dd t,
  \qquad e_1=(1,0).
\]
Show that
\[
  G(x,y)=-\frac1\pi\log|x-y|.
\]

\noindent\textbf{Solution.}
Let $a=|x-y|$.  Since
\[
  p_t(x,y)=\frac1{2\pi t}e^{-a^2/(2t)},
  \qquad
  p_t(0,e_1)=\frac1{2\pi t}e^{-1/(2t)},
\]
we have
\[
  G(x,y)
  =
  \frac1{2\pi}\int_0^\infty
  \left(e^{-a^2/(2t)}-e^{-1/(2t)}\right)t^{-1}\dd t.
\]
Set $u=1/(2t)$.  Then $t^{-1}\dd t=-u^{-1}\dd u$, and
\[
  G(x,y)
  =
  \frac1{2\pi}\int_0^\infty
  \frac{e^{-a^2u}-e^{-u}}{u}\dd u.
\]
The standard Frullani integral gives, for $a>0$,
\[
  \int_0^\infty \frac{e^{-a^2u}-e^{-u}}{u}\dd u
  =
  \log\frac1{a^2}
  =
  -2\log a.
\]
Hence
\[
  G(x,y)=-\frac1\pi\log|x-y|.
\]

\subsection*{Exercise 9.6.E: distributional normalization}

\noindent\textbf{Problem.}
Verify that the constants in the potential kernels are chosen so that
\[
  \frac12\Del_xG(x,0)=-\delta_0
\]
in the distributional sense.

\noindent\textbf{Solution.}
Away from $0$, each kernel is a constant multiple of the radial
harmonic function from Section 9.1:
\[
  |x|^{2-d}\quad(d\ge3),\qquad
  \log|x|\quad(d=2),\qquad
  |x|\quad(d=1).
\]
So $\Del G(x,0)=0$ for $x\ne0$.  It remains to check the mass at the
origin.

For $d=1$, $G(x,0)=-|x|$.  Since $(|x|)'=\operatorname{sgn}(x)$ and
$(|x|)''=2\delta_0$,
\[
  \frac12G''=-\delta_0.
\]

For $d\ge3$, write
\[
  G(x,0)=C_d|x|^{2-d},
  \qquad
  C_d=\frac{\Gamma(d/2-1)}{2\pi^{d/2}}.
\]
Let $\varphi\in C_c^\infty(\R^d)$.  Since $G$ is harmonic off $0$,
Green's identity on $\{|x|>\varepsilon\}$ gives
\[
  \int_{|x|>\varepsilon}G(x,0)\Del\varphi(x)\dd x
  =
  \int_{|x|=\varepsilon}
  \left\{G\,\partial_n\varphi-\varphi\,\partial_nG\right\}\dd S,
\]
where $n$ is the inward normal for the removed ball.  The term
$G\,\partial_n\varphi$ is $O(\varepsilon)$ and vanishes.  Also
\[
  \partial_rG=C_d(2-d)\varepsilon^{1-d}.
\]
With the inward normal convention on the inner boundary,
$\partial_nG=-\partial_rG=C_d(d-2)\varepsilon^{1-d}$.  Hence
\[
  \int_{\R^d}G\Del\varphi
  =
  -C_d(d-2)\omega_d\varphi(0),
\]
where $\omega_d=2\pi^{d/2}/\Gamma(d/2)$ is the surface area of the unit
sphere.  Using
\[
  \Gamma(d/2)=(d/2-1)\Gamma(d/2-1),
\]
we get $C_d(d-2)\omega_d=2$.  Therefore
\[
  \int_{\R^d}\frac12G\Del\varphi=-\varphi(0),
\]
which is exactly $(1/2)\Del G=-\delta_0$.

For $d=2$, $G(x,0)=-(1/\pi)\log|x|$.  The same boundary calculation
uses $\partial_rG=-(1/\pi)r^{-1}$ and the circle length $2\pi r$,
giving
\[
  \int_{\R^2}G\Del\varphi=-2\varphi(0),
  \qquad
  \int_{\R^2}\frac12G\Del\varphi=-\varphi(0).
\]
Thus the normalization is correct in every dimension.

\section*{Section 9.7 Exercises}

\Knowledge{
\path{durrett_ch9_poisson_equation_representation};
\path{durrett_ch9_expected_exit_time_poisson};
\path{durrett_ch9_poisson_boundary_regular}.
}

\subsection*{Exercise 9.7.1}

\noindent\textbf{Question.}
Let $G$ be a bounded open set and let
\[
  \tau=\inf\{t:B_t\notin G\}.
\]
What equation should we solve to find
\[
  w(x)=\E_x\tau^2?
\]

\noindent\textbf{Solution.}
First recall the Section 9.7 result for the first moment.  If
\[
  v(x)=\E_x\tau,
\]
then $v$ solves
\[
  \frac12\Del v=-1\quad\text{in }G,
  \qquad
  v=0\quad\text{on }\partial G
\]
at regular boundary points.

For the second moment, use the Markov property at a small deterministic
time $t$, or equivalently Dynkin's formula.  On the event $\{t<\tau\}$,
\[
  \tau=t+\tau\circ\theta_t,
\]
so the general moment recursion is
\[
  \frac12\Del m_n=-n m_{n-1},
  \qquad
  m_n=0\quad\text{on }\partial G,
\]
where $m_n(x)=\E_x\tau^n$ and $m_0\equiv1$.  Therefore
\[
  w(x)=m_2(x)=\E_x\tau^2
\]
should be found by solving
\[
  \boxed{\frac12\Del w=-2v\quad\text{in }G,\qquad w=0\quad\text{on }\partial G,}
\]
where $v$ is the solution of the first-moment problem above.

Equivalently, eliminating $v$ gives a fourth-order equation.  Since
\[
  v=-\frac14\Del w,
\]
and $(1/2)\Del v=-1$, we get
\[
  \frac18\Del^2 w=-1,
  \qquad\text{or}\qquad
  \Del^2 w=-8
  \quad\text{in }G.
\]
The coupled Poisson system is usually the cleaner way to state and
solve the problem:
\[
  \frac12\Del v=-1,\quad v|_{\partial G}=0,
  \qquad
  \frac12\Del w=-2v,\quad w|_{\partial G}=0.
\]

\section*{Section 9.8 Exercises}

\noindent\textbf{Question.}
The local Chapter 9 PDF has no separately numbered exercise block in
Section 9.8, ``Schrodinger Equation.''  The solutions below work
through the section's main verification problems: the weighted
martingale, the one-dimensional obstruction, the finite-gauge
condition, the Feynman-Kac representation, and the explicit exponential
exit-time formula.

\Knowledge{
\path{durrett_ch9_schrodinger_martingale};
\path{durrett_ch9_schrodinger_nonuniqueness_warning};
\path{durrett_ch9_gauge_function};
\path{durrett_ch9_gauge_theorem_connected_domain};
\path{durrett_ch9_schrodinger_dirichlet_representation};
\path{durrett_ch9_schrodinger_boundary_regular};
\path{durrett_ch9_interval_exponential_exit_moment}.
}

\subsection*{Exercise 9.8.A: the Schrodinger martingale}

\noindent\textbf{Problem.}
Let $G$ be open, let $\tau=\inf\{t:B_t\notin G\}$, and suppose
$u\in C^2(G)$ satisfies
\[
  \frac12\Del u+c(x)u=0\quad\text{in }G.
\]
Show that
\[
  M_t=u(B_t)\exp\left(\int_0^t c(B_s)\dd s\right),
  \qquad 0\le t<\tau,
\]
is a local martingale.

\noindent\textbf{Solution.}
Put
\[
  A_t=\int_0^t c(B_s)\dd s.
\]
Since $A_t$ has finite variation, the product rule and It\^o's formula
give, for $t<\tau$,
\[
\begin{aligned}
  \dd\{u(B_t)e^{A_t}\}
  &=
  e^{A_t}\nabla u(B_t)\cdot\dd B_t
  +
  e^{A_t}\frac12\Del u(B_t)\dd t
  +
  u(B_t)e^{A_t}c(B_t)\dd t \\
  &=
  e^{A_t}\nabla u(B_t)\cdot\dd B_t
  +
  e^{A_t}\left(\frac12\Del u+c u\right)(B_t)\dd t.
\end{aligned}
\]
The drift vanishes by the PDE.  After stopping when
\[
  \int_0^t e^{2A_s}|\nabla u(B_s)|^2\dd s
\]
gets large or when $B_t$ approaches $\partial G$, the stochastic
integral is a martingale.  Hence $M_t$ is a local martingale on
$[0,\tau)$.

\subsection*{Exercise 9.8.B: why boundedness alone is not enough}

\noindent\textbf{Problem.}
Solve the boundary-value problem
\[
  \frac12u''+\gamma u=0\quad\text{on }(-a,a),
  \qquad
  u(-a)=u(a)=1,
\]
and explain why this warns us not to assert blindly that
\[
  u(x)=\E_x e^{\gamma\tau},
  \qquad
  \tau=\inf\{t:B_t\notin(-a,a)\}.
\]

\noindent\textbf{Solution.}
Let $b=\sqrt{2\gamma}$.  The general solution is
\[
  u(x)=A\cos(bx)+B\sin(bx).
\]
The boundary equations are
\[
  A\cos(ba)+B\sin(ba)=1,
  \qquad
  A\cos(ba)-B\sin(ba)=1.
\]
Thus $B\sin(ba)=0$ and $A\cos(ba)=1$.  When $\cos(ba)\ne0$, the even
solution is
\[
  u(x)=\frac{\cos(bx)}{\cos(ba)}.
\]
When $\cos(ba)=0$, no solution with boundary value $1$ exists.

The probabilistic quantity $\E_x e^{\gamma\tau}$ is nonnegative.  But
if $ba>\pi/2$ and $\cos(ba)\ne0$, the ODE solution above takes negative
values for some $x$, so it cannot equal $\E_xe^{\gamma\tau}$.  The
missing hypothesis is integrability of the exponential weight, encoded
by the gauge
\[
  w(x)=\E_x e^{\gamma\tau}.
\]
For large positive $\gamma$, this gauge is infinite.

\subsection*{Exercise 9.8.C: the gauge and finite-gauge uniqueness}

\noindent\textbf{Problem.}
Let
\[
  w(x)=\E_x\exp\left(\int_0^\tau c(B_s)\dd s\right).
\]
Assume $G$ is bounded and connected, $f,c$ are bounded and continuous,
and $w\not\equiv\infty$.  Explain why any bounded solution of
\[
  \frac12\Del u+c u=0\quad\text{in }G,\qquad u=f\quad\text{on }\partial G
\]
must be
\[
  u(x)=
  \E_x\left[
    f(B_\tau)\exp\left(\int_0^\tau c(B_s)\dd s\right)
  \right].
\]

\noindent\textbf{Solution.}
By the gauge theorem in the section, $w\not\equiv\infty$ on a connected
bounded domain implies
\[
  \sup_{x\in G}w(x)<\infty.
\]
Let
\[
  A_t=\int_0^t c(B_s)\dd s.
\]
Exercise 9.8.A shows that $u(B_t)e^{A_t}$ is a local martingale before
exit.  Stop at $t\wedge\tau$ and use boundedness to get
\[
  u(x)
  =
  \E_x\{f(B_\tau)e^{A_\tau};\tau\le t\}
  +
  \E_x\{u(B_t)e^{A_t};\tau>t\}.
\]
The first term converges by dominated convergence to
\[
  \E_x\{f(B_\tau)e^{A_\tau}\},
\]
because $|f(B_\tau)e^{A_\tau}|\le \|f\|_\infty e^{A_\tau}$ and
$\E_xe^{A_\tau}=w(x)<\infty$.

For the second term, condition on $\mathcal{F}_t$ and use the Markov
property:
\[
  \E_x(e^{A_\tau}\mid\mathcal{F}_t)
  =
  e^{A_t}w(B_t)
  \quad\text{on }\{\tau>t\}.
\]
Since the finite-gauge theorem gives $w$ bounded and bounded away from
$0$ on $G$ in the way used in Durrett's proof, the tail term is
controlled by
\[
  C\,\E_x\{e^{A_\tau};\tau>t\}\to0
\]
again by dominated convergence.  Therefore
\[
  u(x)=
  \E_x\left[
    f(B_\tau)\exp\left(\int_0^\tau c(B_s)\dd s\right)
  \right],
\]
and the bounded solution is unique.

\subsection*{Exercise 9.8.D: the smooth representation solves the PDE}

\noindent\textbf{Problem.}
Let
\[
  v(x)=
  \E_x\left[
    f(B_\tau)\exp\left(\int_0^\tau c(B_s)\dd s\right)
  \right].
\]
Assume $v\in C^2(G)$.  Show that
\[
  \frac12\Del v+c v=0\quad\text{in }G.
\]

\noindent\textbf{Solution.}
Let $B(x,r)\Subset G$ and let
\[
  \sigma=\inf\{t:B_t\notin B(x,r)\}.
\]
By the strong Markov property,
\[
  v(x)=
  \E_x\left[
    e^{\int_0^\sigma c(B_s)\dd s}v(B_\sigma)
  \right].
\]
For small $r$, $\E_x\sigma=O(r^2)$ and
\[
  e^{\int_0^\sigma c(B_s)\dd s}
  =
  1+c(x)\sigma+o(r^2)
\]
in expectation.  Taylor expansion gives
\[
  \E_x v(B_\sigma)-v(x)
  =
  \frac12\Del v(x)\E_x\sigma+o(r^2),
\]
because the first moments vanish and the cross terms cancel by
symmetry.  Combining these expansions,
\[
  0=
  \left(\frac12\Del v(x)+c(x)v(x)\right)\E_x\sigma+o(r^2).
\]
Since $\E_x\sigma$ is of order $r^2$, divide by $\E_x\sigma$ and let
$r\downarrow0$:
\[
  \frac12\Del v(x)+c(x)v(x)=0.
\]

\subsection*{Exercise 9.8.E: boundary condition at regular points}

\noindent\textbf{Problem.}
Assume the finite-gauge hypotheses and let $y$ be a regular boundary
point.  Show that
\[
  v(x)\to f(y)
  \qquad\text{as }x\to y,\ x\in G.
\]

\noindent\textbf{Solution.}
Let $x_n\to y$ from inside $G$.  Regularity gives, for every
$\delta>0$,
\[
  \Prob_{x_n}(\tau\le\delta)\to1,
  \qquad
  \Prob_{x_n}(B_\tau\in B(y,\delta))\to1.
\]
On $\{\tau\le\delta,\ B_\tau\in B(y,\delta)\}$, continuity of $f$ and
boundedness of $c$ imply
\[
  e^{\int_0^\tau c(B_s)\dd s}f(B_\tau)\to f(y)
\]
as $\delta\downarrow0$.  The part with $\tau>1$ is negligible because
the finite-gauge theorem gives $\sup_G w<\infty$ and
\[
  \E_{x_n}\{e^{\int_0^\tau c(B_s)\dd s}|f(B_\tau)|;\tau>1\}
  \le
  C\|f\|_\infty \Prob_{x_n}(\tau>1)\to0.
\]
The remaining intermediate interval is handled by first taking
$\delta$ small and then $n$ large.  Hence $v(x_n)\to f(y)$.

\subsection*{Exercise 9.8.F: exponential exit time from an interval}

\noindent\textbf{Problem.}
Let $\tau=\inf\{t:B_t\notin(-a,a)\}$ and $\gamma>0$.  Show that
\[
  \E_x e^{\gamma\tau}
  =
  \frac{\cos(x\sqrt{2\gamma})}{\cos(a\sqrt{2\gamma})}
  \quad\text{if }0<\gamma<\frac{\pi^2}{8a^2},
\]
and that $\E_xe^{\gamma\tau}=\infty$ if
$\gamma\ge \pi^2/(8a^2)$.

\noindent\textbf{Solution.}
Let $b=\sqrt{2\gamma}$.  If $0<ab<\pi/2$, then
\[
  u(x)=\frac{\cos(bx)}{\cos(ba)}
\]
is positive on $[-a,a]$ and solves
\[
  \frac12u''+\gamma u=0,
  \qquad
  u(-a)=u(a)=1.
\]
The finite-gauge theorem therefore identifies the Feynman-Kac
representation:
\[
  \E_xe^{\gamma\tau}=u(x)
  =
  \frac{\cos(x\sqrt{2\gamma})}{\cos(a\sqrt{2\gamma})}.
\]

Now suppose $\gamma\ge\pi^2/(8a^2)$.  This is the first Dirichlet
eigenvalue threshold for the interval.  If the gauge were finite, the
finite-gauge theorem would produce a bounded nonnegative solution of
\[
  \frac12u''+\gamma u=0,\qquad u(\pm a)=1.
\]
But the explicit ODE solution from Exercise 9.8.B either does not exist
at the threshold values where $\cos(ba)=0$, or changes sign when
$ab>\pi/2$.  Neither possibility can equal the nonnegative function
$\E_xe^{\gamma\tau}$.  Hence the gauge cannot be finite, and
\[
  \E_xe^{\gamma\tau}=\infty
  \qquad\text{for all }x\in(-a,a).
\]

\end{document}
